\documentclass[10pt,a4paper]{amsart}
\usepackage[margin=19mm]{geometry}
\usepackage{amsmath,amssymb,mathtools}
\usepackage[scr=rsfs,cal=euler]{mathalfa}
\usepackage{xfrac}
\usepackage[unicode,hidelinks,bookmarks=false]{hyperref}
\newcommand{\ie}{i.e.\ }
\newcommand{\cf}{cf.\ }
\newcommand{\resp}{respectively\ }
\newcommand{\Subsection}{\subsection}
\providecommand{\nobreakdash}{\nobreak\hbox{-}}
\setlength{\emergencystretch}{2em}
\multlinegap0pt
% ---------------------------------------------------------------------------
% Editorial AMS wrapper prelude. The theorem environments and the mathematical macros below are
% transcribed from the paper's own preamble (cached-originals/source0.tex lines 16--88); the AMS amsart
% class, the named format packages of the pinned TeX Live runtime and this editorial formatting are not
% part of the author's text. No mathematics is changed.
% ---------------------------------------------------------------------------
\newtheorem{theorem}{Theorem}[section]
\newtheorem{corollary}[theorem]{Corollary}
\newtheorem{proposition}[theorem]{Proposition}
\newtheorem{lemma}[theorem]{Lemma}
\newtheorem{construction}[theorem]{Construction}
\theoremstyle{definition}
\newtheorem{definition}[theorem]{Definition}
\newtheorem{example}[theorem]{Example}
\newtheorem{remark}[theorem]{Remark}
\newtheorem{problem}[theorem]{Problem}
\newtheorem{observation}[theorem]{Observation}
\newtheorem{claim}[theorem]{Claim}
\newtheorem{conjecture}[theorem]{Conjecture}

\def\R{\mathbb R}
\DeclareMathOperator{\Diag}{Diag}
\DeclareMathOperator{\vol}{vol}
\newcommand{\Prob}[1]{\mathbb{P}\left(#1\right)}
\newcommand{\CProb}[2]{\mathbb{P}\left(\left.#1\right|#2\right)}
\newcommand{\Exp}[1]{\mathbb{E}\left[#1\right]}
\newcommand{\cartesianproduct}{\,\square\,}


\newcommand{\marrow}{\marginpar[\hfill$\longrightarrow$]{$\longleftarrow$}}
\newcommand{\niceremark}[3]{\textcolor{blue}{\textsc{#1 #2:} \marrow\textsf{#3}}}
\newcommand{\niceremarkcolor}[4]{\textcolor{#4}{\textsc{#1 #2:} \marrow\textsf{#3}}}
\newcommand{\aida}[2][says]{\niceremarkcolor{Aida}{#1}{#2}{Red}}
\newcommand{\thijs}[2][says]{\niceremarkcolor{Thijs}{#1}{#2}{Purple}}
\newcommand{\harper}[2][says]{\niceremarkcolor{Harper}{#1}{#2}{Blue}}
\newcommand{\emanuel}[2][says]{\niceremarkcolor{Emanuel}{#1}{#2}{DarkDandelion}}
\newcommand{\robin}[2][says]{\niceremarkcolor{Robin}{#1}{#2}{Brown}}
\newcommand{\jim}[2][says]{\niceremarkcolor{Jim}{#1}{#2}{Green}}
\newcommand{\nils}[2][says]{\niceremarkcolor{Nils}{#1}{#2}{RubineRed}}


\newcommand{\todo}[2][says]{\niceremarkcolor{To do}{#1}{#2}{CadetBlue}}




\begin{document}
\title{The edge-isoperimetric number of graphs and their powers}
\author{Aida Abiad, Nils van de Berg, Emanuel Juliano, Harper Reijnders, Robin Simoens, Thijs van Veluw and Jim Wittebol}
\date{Source: arXiv:2601.17519v2; distributed under CC BY 4.0}
\maketitle
\noindent\emph{Source, version and licence.} \emph{The edge-isoperimetric number of graphs and their powers}, by Aida Abiad, Nils van de Berg, Emanuel Juliano, Harper Reijnders, Robin Simoens, Thijs van Veluw and Jim Wittebol. The source of this editable unit is the e-print of \textbf{version v2} of arXiv:2601.17519, https://arxiv.org/abs/2601.17519v2, whose material is distributed under the Creative Commons Attribution 4.0 International licence, http://creativecommons.org/licenses/by/4.0/, as recorded in the arXiv licence metadata for this paper and shown on that abstract page. The whole of the body below is version v2, the newest version of the paper. Full rights notice: licenses/arxiv-2601.17519-CC-BY-4.0.txt.\par\smallskip
\noindent\textit{Editorial scope.} Counted unit: original Theorem 3.7 of the article, the statement carried by the source environment \texttt{theorem} with source label \texttt{theorem:lowerboundi(G\textasciicircum 2)closed}, cached-originals/source0.tex lines 357--367, printed as \textbf{Theorem 3.7} exactly as in the article. It is delivered together with the authors' own complete proof as the single primary proof body, source0.tex lines 370--428 (59 lines): the \texttt{proof} environment that immediately follows the statement and optimises the polynomial case by case. No proof marker is added and no mathematics is written, repaired or re-derived; the authors' notation and the printed slips of the source are kept literal. Necessary complete context is retained uncounted: source0.tex lines 305--308, Theorem 3.4, the general polynomial lower bound whose optimisation the counted theorem carries out for $t=2$; the closure of the counted proof over references was computed and the further passages it needs are delivered with it. The article's own numbering is reproduced by explicit counter settings: the counter \texttt{theorem}, declared by \texttt{\textbackslash newtheorem\{theorem\}\{Theorem\}[section]} with corollary, proposition, lemma, construction, definition, example, remark, problem, observation, claim and conjecture sharing it, is set before each retained passage, so the counted statement prints as Theorem 3.7 and the general theorem as Theorem 3.4, exactly as in the article. Every \texttt{\textbackslash ref} inside every retained passage resolves inside this delivered file, and the single citation that occurs in the retained passages, \texttt{spectra}, is delivered as the authors' own bibliography entry transcribed verbatim from the article's own \texttt{thebibliography} with the authors' own printed number. The retained passages contain no figure environment and no graphics-inclusion command, so no artwork is redistributed. The source's own class is the standard \texttt{article} class; the e-print ships no class or style file, so no publisher class is shipped, used or redistributed, and the wrapper is the AMS \texttt{amsart} class with the article's theorem declarations transcribed from its preamble. Every declared body of this unit is an exact, unmodified span of source0.tex: no body is a bounded passage comparison, \texttt{omit} was not used anywhere, and consequently no fidelity receipt is asserted in delivery-evidence.json. The complete concatenated original source is bundled as sources/193320/source0.tex. Omitted from this editable unit and therefore absent from it are source0.tex lines 1--304; 309--356; 368--369; 429--1333; 1336--1644, that is the article's own preamble, title and author block, the introduction, the general spectral theory sections, the upper-bound counterparts and the remaining applications, and the remaining bibliography entries, all of which remain present in the bundled source but are not delivered here. The AMS \texttt{amsart} wrapper, the named format packages of the pinned TeX Live runtime and this editorial source, licence and scope paragraph are editorial formatting only.\par\smallskip
\setcounter{section}{3}
\setcounter{equation}{0}
\setcounter{theorem}{3}
\begin{theorem}\label{theorem:lowerboundi(G^t)}
    Let $G$ be a regular graph on $n\ge 2$ vertices with adjacency eigenvalues $\lambda_1 \ge \dots \ge \lambda_n$. Let $p\in \R_t[x]$, and define $W(p) = \max_{1\le i < j \le n} (p(A))_{ij}$ the maximum off-diagonal entry of $p(A)$. If $W(p)> 0$ and $p(\lambda_1)\ge \Lambda(p)$, then
    $$i(G^t) \ge \frac{p(\lambda_1) - \Lambda(p)}{2W(p)}.$$
\end{theorem}

\setcounter{section}{3}
\setcounter{equation}{0}
\setcounter{theorem}{6}
\begin{theorem}\label{theorem:lowerboundi(G^2)closed}
    Let $G$ be a non-complete regular connected graph with adjacency eigenvalues $\lambda_1\ge \dots \ge \lambda_n$. Let $\Lambda$ be the maximum number of common neighbours of adjacent vertices of $G$, and let $M$ be the maximum number of neighbours of non-adjacent vertices of $G$. The following choices of $p$ are optimal for the bound of Theorem \ref{theorem:lowerboundi(G^t)} when $t=2$.
    \begin{enumerate}
        \item[i.] If $\Lambda-\lambda_n \le M+\lambda_2$, then $p=x^2+(M-\Lambda)x$, giving
        $$i(G^2) \ge \frac{(\lambda_1-\lambda_2)(\lambda_1+\lambda_2+M-\Lambda)}{2M}.$$
        \item[ii.] If $M+\lambda_2 \le \Lambda - \lambda_n \le \lambda_1$, then $p=x^2+(M-\Lambda)x$, giving
        $$i(G^2)\ge \frac{(\lambda_1-\lambda_n)(\lambda_1+\lambda_n+M-\Lambda)}{2M}.$$
        \item[iii.] If $\max(\lambda_1,M+\lambda_2) \le \Lambda-\lambda_n$, then $p=x^2-(\lambda_2+\lambda_n)x$, giving
        $$i(G^2)\ge \frac{(\lambda_1-\lambda_2)(\lambda_1-\lambda_n)}{2(\Lambda-\lambda_2-\lambda_n)}.$$
    \end{enumerate}
\end{theorem}

\setcounter{section}{3}
\setcounter{equation}{0}
\setcounter{theorem}{7}
\begin{proof}
    Let $p=ax^2+bx+c \in \R_2[x]$. Let $B(p)=\frac{p(\lambda_1) - \Lambda(p)}{2W(p)}$ denote the bound from Theorem \ref{theorem:lowerboundi(G^t)}. We optimize $B(p)$ for the parameters $a$, $b$, and $c$, provided that $p(\lambda_1)\ge \Lambda(p)$ and $W(p)>0$. First of all, note that scaling by a positive number and shifting does not change the value of $B(p)$, so that we may assume without loss of generality that $c=0$ and $a\in \{-1,0,1\}$. We make a case distinction according to the value of $a$.
    \begin{itemize}
    \item If $a=0$, then by $W(p)>0$ and scaling, we may assume that $b=1$, and so $p=x$ and $B(p)=(\lambda_1-\lambda_n)/2$.
    \item If $a=-1$, then $p=-x^2+bx$. We show that this case gives bounds that are worse than the bound from the case $a=0$. To show this, we investigate the numerator and denominator of $B(p)$.
    
    For the numerator of $B(p)$ we note the following. For any $1\le i\le j\le n$, we have that $p(\lambda_i)\ge p(\lambda_j)$ if and only if $b\ge \lambda_i+\lambda_j$ or $\lambda_i=\lambda_j$. Since $G$ is connected, $\lambda_1>\lambda_2$ by the Perron-Frobenius Theorem \cite{spectra}. Thus, in order for $p(\lambda_1)\ge \Lambda(p)$, by the above we need $b\ge \lambda_1+\lambda_2$, and consequently $\Lambda(p)=p(\lambda_2)$, again by the above. Now $p(\lambda_1)-\Lambda(p)=p(\lambda_1)-p(\lambda_2)=(\lambda_1-\lambda_2)(b-\lambda_1-\lambda_2)$.

    For the denominator of $B(p)$, we have the following. Write $\eta$ for the minimum number of common neighbours of two adjacent vertices of $G$. In order to have $W(p)>0$ we need $b>\eta$ because $A^2$ only has non-negative entries. In this case $W(p)=b-\eta$.
    So, for $p(\lambda_1)\ge \Lambda(p)$ and $W(p)>0$, we need that $b\ge\lambda_1+\lambda_2$ and $b>\eta$, in which case
    $$B(p)=\frac{(\lambda_1-\lambda_2)(b-\lambda_1-\lambda_2)}{2(b-\eta)}=\frac{\lambda_1-\lambda_2}2 \left (1+\frac{\eta-\lambda_1-\lambda_2}{b-\eta} \right).$$
    Note that $\eta < \lambda_1-1 < \lambda_1+\lambda_2$ since $G$ is not complete, so $B(p)< (\lambda_1-\lambda_2)/2$, implying that the bounds coming from polynomials in this case are worse than the bound from the polynomial $p=x$ as in the above case $a=0$.
    
    \item If $a=1$, then $p=x^2+bx$. We again investigate the numerator and denominator of $B(p)$.

    For the numerator of $B(p)$, again if $1\le i \le j \le n$, then $p(\lambda_i)\ge p(\lambda_j)$ if and only if $b\ge -\lambda_i-\lambda_j$ or $\lambda_i=\lambda_j$. In order to have $p(\lambda_1)\ge \Lambda(p)$, we therefore need that $b\ge -\lambda_1-\lambda_n$. Furthermore, by convexity of the parabola $p=x^2+bx$, we have $\Lambda(p)\in \left\{p(\lambda_2),p(\lambda_n)\right\}$, and if $b\ge -\lambda_2-\lambda_n$, then $\Lambda(p)=p(\lambda_2)$. If instead $-\lambda_1-\lambda_n \le b \le -\lambda_2-\lambda_n$, then $\Lambda(p)=p(\lambda_n)$. Summarizing, we have 
    \[p(\lambda_1)-\Lambda(p)=\begin{cases}
        $negative$ &$if $b<-\lambda_1-\lambda_n,\\
        (\lambda_1-\lambda_n)(b+\lambda_1+\lambda_n) &$if $-\lambda_1-\lambda_n \le b \le -\lambda_2 - \lambda_n,\\
        (\lambda_1-\lambda_2)(b+\lambda_1+\lambda_2) &$if $-\lambda_2-\lambda_n \le b.
    \end{cases}\]

    For the denominator of $B(p)$, note that $W(p)=\max(\Lambda+b,M)$. Since $G$ is non-complete and connected, $M>0$, so $W(p)>0$ is automatically satisfied. We now have
    \[W(p)=\begin{cases}
        M &$if $b\le M-\Lambda,\\
        \Lambda+b &$if $M-\Lambda\le b.
    \end{cases}\]

    To optimize the value of $B(p)$, we look at the behaviour of $B(p)$ on three (possibly empty) key intervals. Recall that $b \ge -\lambda_1-\lambda_n$.
    \begin{align*}
        I_1 &= \Big[-\lambda_1-\lambda_n,\ M-\Lambda\Big ],\\
        I_2 &= \Big [\max\big(-\lambda_1-\lambda_n,\ M-\Lambda\big),\ -\lambda_2-\lambda_n\Big ],\\
        I_3 &= \Big [ \max\big(-\lambda_2-
        \lambda_n,\ M-\Lambda\big),\ \infty \Big ).
    \end{align*}
    \begin{enumerate}
    \item For $b\in I_1$, the denominator of $B(p)$ is constant and equal to $2M$. The numerator is increasing in $b$, so $B(p)$ is increasing on $I_1$.
    \item If $b\in I_2$, then
    $$B(p)=\frac{(\lambda_1-\lambda_n)(b+\lambda_1+\lambda_n)}{2(b+\Lambda)}=\frac{\lambda_1-\lambda_n}{2}\left (1+\frac{\lambda_1+\lambda_n-\Lambda}{b+\Lambda}\right ).$$
    So $B(p)$ is strictly increasing on $I_2$ if $\Lambda>\lambda_1+\lambda_n$, constant if $\Lambda=\lambda_1+\lambda_n$, and strictly decreasing if $\Lambda< \lambda_1+\lambda_n$.
    \item If $b\in I_3$, then
    $$B(p)=\frac{(\lambda_1-\lambda_2)(b+\lambda_1+\lambda_2)}{2(b+\Lambda)}=\frac{\lambda_1-\lambda_2}2 \left (1+\frac{\lambda_1+\lambda_2-\Lambda}{b+\Lambda} \right).$$
    By regularity and non-completeness $\Lambda\le \lambda_1-1<\lambda_1+\lambda_2$, so that $B(p)$ is decreasing on $I_3$.
    \end{enumerate}
    To complete the proof, we now make the following case distinction.
    \begin{itemize}
        \item If $-\lambda_2-\lambda_n \le M-\Lambda$, then $I_2$ is empty. Since $B(p)$ is increasing on $I_1=[-\lambda_1-\lambda_n,\ M-\Lambda]$ and decreasing on $I_3=[M-\Lambda,\ \infty)$, the optimal value for $b$ is $M-\Lambda$, giving the bound
        $$i(G^2)\ge B(p) = \frac{(\lambda_1-\lambda_2)(\lambda_1+\lambda_2+M-\Lambda)}{2M},$$
        which is greater than $B(x)$ because $\Lambda <\lambda_1+\lambda_2$.
            \item If $M-\Lambda\le -\lambda_2-\lambda_n$ and $\lambda_1+\lambda_n \le \Lambda$, then $B(p)$ is increasing on $I_1\cup I_2=[-\lambda_1-\lambda_n,\ -\lambda_2-\lambda_n]$, and decreasing on $I_3=[-\lambda_2-\lambda_n,\ \infty)$. It follows that the optimal value of $b$ is $-\lambda_2-\lambda_n$, giving the bound
            $$i(G^2)\ge B(p) = \frac{(\lambda_1-\lambda_2)(\lambda_1-\lambda_n)}{2(\Lambda-\lambda_2-\lambda_n)},$$
            which is again greater than $B(x)$ because $\Lambda<\lambda_1+\lambda_2$.
            \item Lastly, if $M-\Lambda \le -\lambda_2-\lambda_n$ and $\Lambda \le \lambda_1+\lambda_n$, then we must have $-\lambda_1-\lambda_n \le M-\Lambda$, and so $B(p)$ is increasing on $I_1=[-\lambda_1-\lambda_n,\ M-\Lambda]$, and decreasing on $I_2\cup I_3=[M-\Lambda,\ \infty)$. Now the optimal value of $b$ is $M-\Lambda$, giving the bound
            $$i(G^2)\ge B(p)=\frac{(\lambda_1-\lambda_n)(\lambda_1+\lambda_n+M-\Lambda)}{2M}.$$
            Again, this is greater than $B(x)$, because $\Lambda\le \lambda_1+\lambda_n$ and $\lambda_1-\lambda_2 \le \lambda_1-\lambda_n$.
    \end{itemize}
    \end{itemize}
    This concludes the proof.
\end{proof}

\begin{thebibliography}{99}
\bibitem[35]{spectra} A.E. Brouwer, W.H. Haemers, \emph{Spectra of graphs}, Universitext, Springer, New York, 2012. 
\end{thebibliography}
\end{document}
